\begin{table}[H]\centering
\caption{Descriptive statistics}
\label{tab:desc_stats}
\begin{tabular}{lcccccc}\hline\hline
& \multicolumn{3}{c}{Quality (ISO 9001)} & \multicolumn{3}{c}{Environment (ISO 14001)} \\
\cmidrule(lr){2-4} \cmidrule(lr){5-7}
& 1993 & 1999 & 2017 & 1993 & 1999 & 2017 \\\hline
Countries Number & 27 & 27 & 27 & & 27 & 27 \\
Total & 8,866 & 115,052 & 315,987 & & 4,996 & 84,316 \\
Mean & 328.4 & 4,261.2 & 11,703.2 & & 185.0 & 3,122.8 \\
Std.\ Dev. & 508.9 & 7,289.6 & 21,695.1 & & 268.9 & 4,020.2 \\
Min--Max & 0--1,586 & 39--30,150 & 191--97,646 & & 0--962 & 35--14,571 \\
Cert./GDP p.c. & 11.52 & 158.51 & 441.67 & & 6.44 & 128.86 \\
\hline\hline
\multicolumn{7}{p{0.92\textwidth}}{\footnotesize \textit{Notes:} This table summarises the distribution of certifications across the 27 EU countries in the three selected years. The first row reports the number of countries with non-missing data in each year, which is the basis for all subsequent statistics. The second row reports the total number of certifications across all countries, while the third and fourth rows report the mean and standard deviation of certifications per country, respectively. The fifth row reports the minimum and maximum number of certifications across countries. Finally, the sixth row reports the mean certification intensity, defined as the number of certifications per thousand dollars of GDP per capita, averaged across countries. Environment 1993 is empty as its adoption started in 1997. Sources: ISO Survey, World Bank.} \\
\end{tabular}
\end{table}
